\begin{lemma}[Initial reserved pairs in the PE case]
Let $\mathcal E$ be a positive-rung ETT with rung number $n$, terminal
ladder tree $T_n$, generating coloring $\varphi_n$, and terminal extension
kind PE.  Fix a \noderef{HierarchyBase} $B$ and a
\noderef{TashkinovHierarchy} $H$ over $B$, and put
\[
 M=\noderef{TreeMissingColors}(\varphi_n,T_n),\qquad
 D=D^H_{n,0}=\noderef{HierarchyStoredColorsAt}(H,0),
 \qquad C=M\cap\overline\varphi_n(B.\mathit{closure}).
\]
Assume $|M|\geq2n+11$.  Suppose that some
$\alpha\in\overline\varphi_n(B.\mathit{core})$ is not closed on
$B.\mathit{core}$, that every color in $C$ is closed on
$B.\mathit{core}$, and that $|C|\geq4$.

Then there is a total family of finite color sets $\Gamma_\eta$, constrained
at every $\eta\in D$, such that
\begin{enumerate}
\item $|\Gamma_\eta|=2$ and $\Gamma_\eta\subseteq M$ for every
  $\eta\in D$;
\item $\Gamma_\eta\cap\Gamma_\theta=\varnothing$ for distinct
  $\eta,\theta\in D$;
\item $B.\mathit{core}$ is not
  $\bigl(\bigcup_{\eta\in D}\Gamma_\eta\bigr)^-$-closed under
  $\varphi_n$; and
\item at least four common missing colors remain unreserved:
  \[
  \left|C-\bigcup_{\eta\in D}\Gamma_\eta\right|\geq4.
  \]
\end{enumerate}
The family is total only to match the Lean representation; the conclusions
impose conditions exactly at the paper indices $\eta\in D$.
\end{lemma}
\begin{proof}
Choose $\alpha$ as in the hypotheses.  Since every color of $C$ is closed
on the core while $\alpha$ is not, we have
\[
                              \alpha\notin C.                 \tag{1}
\]
The stagewise bound \noderef{HierarchyStoredColorsCardinality} gives
$|D|\leq n$.  Write $d=|D|$ and $A=M-\{\alpha\}$.  Deleting one element
reduces cardinality by at most one, so
\[
                  |A|\geq |M|-1\geq2n+10\geq2d+10.           \tag{2}
\]

We choose a set $R\subseteq A$ of exactly $2d$ colors, giving priority to
colors outside $C$.  More explicitly, if $|A-C|\geq2d$, take any
$2d$-subset of $A-C$.  Otherwise take every color of $A-C$ and add any
$2d-|A-C|$ colors from $A\cap C$.  The latter choice is possible because
the disjoint union
$A=(A-C)\mathbin{\dot\cup}(A\cap C)$ and (2) give
\[
 |A\cap C|=|A|-|A-C|\geq2d-|A-C|.
\]
Thus in both cases $R\subseteq A$ and $|R|=2d$.

The same choice leaves at least four colors of $C$ outside $R$.  In the
first case $R\cap C=\varnothing$, so this follows from $|C|\geq4$.  In the
second case, (1) gives $C\subseteq A$, and hence
\[
 |C-R|=|A|-2d\geq10                                             \tag{3}
\]
by (2).  In particular $|C-R|\geq4$ in either case.

Enumerate the $d$ elements of the finite set $D$ and the $2d$ elements of
$R$.  Assign to the $i$-th member of $D$ the two colors in positions
$2i$ and $2i+1$ of the enumeration of $R$.  This defines
$\Gamma_\eta$ for $\eta\in D$; put $\Gamma_\eta=\varnothing$ outside
$D$.  Every indexed set has cardinality two, two different indexed sets
are disjoint, and their union is exactly $R$.  Since $R\subseteq A\subseteq
M$, items 1 and 2 follow, while (3) is item 4.

Finally, $\alpha\notin R$ because $R\subseteq A$.  The chosen color
$\alpha$ is missing on the core and has a nonempty exact
$\alpha$-colored boundary there.  Hence it belongs to
$\overline\varphi_n(B.\mathit{core})-R$ but is not closed.  Expanding
\noderef{TreeColorSetMinusClosed} shows that the core is not
$R^-$-closed.  Since $R=\bigcup_{\eta\in D}\Gamma_\eta$, this is item 3
and completes the proof.
\end{proof}
